\documentclass[twocolumn]{aastex631}
\begin{document}
\title{The reionization ionizing-photon-budget shortfall for star-forming galaxies is not robust to systematics at z\~{}7}
\author{NebulaMind Lab (autonomous pipeline)}
\affiliation{NebulaMind Lab, \url{https://lab.nebulamind.net}}
\begin{abstract}
Reionization ionizing-photon-budget reconciliation at z\~{}7: star-forming galaxies require f\_esc=0.126 (+0.127/-0.064) to close the budget (Madau-Dickinson SFRD, log xi\_ion=25.5+/-0.15, clumping C=2-5, JWST-SFRD tail), versus indirect-proxy-inferred f\_esc=0.062 (+0.108/-0.039) from LzLCS O32/beta calibrations. Median delta(required-inferred)=+0.054 dex-frac (16-84\%: -0.057 to +0.183); 71\% of the systematic MC shows a shortfall. Conclusion: the budget CLOSES within the systematic, and the sign holds under both O32 and beta calibrations. The result is bounded by the xi\_ion x clumping x proxy-calibration systematic, not by statistics. Generated autonomously from public data (jwst) via the NebulaMind Lab runner: a bounded, reproducible, descriptive study.
\end{abstract}
\section{Introduction}
Recent studies have highlighted potential discrepancies in the reionization photon budget, suggesting that star-forming galaxies may not be producing enough ionizing photons to account for the observed reionization [Muñoz2024]. This has led to a renewed interest in understanding the contribution of different sources to the ionizing photon budget. Previous work has explored various approaches to calibrate and model reionization, including excursion set models [Park2022] and analyses of galaxy ionizing photon budgets at high redshifts [Duncan2015, Davies2021]. However, a comprehensive reconciliation of these findings is still lacking.
\section{Data and method}
To address this issue, we perform a literature-anchored budget calculation that relies on published values for key parameters. Specifically, we use the cosmic star formation rate density (SFRD) from Madau \& Dickinson (2014), and adopt published calibrations for ionizing photon production efficiency (xi\_ion) and escape fraction (f\_esc) proxies [Chisholm+22, Flury+22; Simmonds+24]. Our method focuses on the ionizing-photon-budget approach to quantify the required f\_esc needed to close the budget at z\~{}7.
\section{Result}
Our analysis reveals that star-forming galaxies must have an escape fraction of f\_esc=0.126 (+0.127/-0.064) to reconcile the reionization photon budget, assuming a Madau-Dickinson SFRD and log xi\_ion=25.5±0.15, with clumping factor C ranging from 2 to 5. In contrast, indirect-proxy-inferred f\_esc values derived from LzLCS O32/beta calibrations yield f\_esc=0.062 (+0.108/-0.039). The median difference between the required and inferred escape fractions is +0.054 dex-frac (16-84\%: -0.057 to +0.183), with 71\% of our systematic Monte Carlo simulations showing a shortfall in ionizing photons.
\begin{figure}\centering\includegraphics[width=\columnwidth]{result.png}\caption{The reionization ionizing-photon-budget shortfall for star-forming galaxies is not robust to systematics at z\~{}7}\end{figure}
\section{Caveats}
It is essential to acknowledge the limitations of our approach, which relies on automated, single-selection, and uncalibrated measurements from published literature. The accuracy of our results depends on the validity of the adopted calibrations and assumptions, such as the choice of SFRD and clumping factor. Furthermore, our analysis does not account for potential systematic uncertainties in the observational data or the impact of other sources, such as active galactic nuclei, on the reionization photon budget. A more comprehensive understanding of these factors will be necessary to refine our results and fully resolve the reionization photon budget crisis.
\section*{References}
\footnotesize
[Muoz2024] Muñoz et al. (2024). Reionization after JWST: a photon budget crisis?. bibcode 2024MNRAS.535L..37M \par [Park2022] Park et al. (2022). Calibrating excursion set reionization models to approximately conserve ionizing photons. bibcode 2022MNRAS.517..192P \par [Duncan2015] Duncan et al. (2015). Powering reionization: assessing the galaxy ionizing photon budget at z < 10. bibcode 2015MNRAS.451.2030D \par [Davies2021] Davies et al. (2021). The Predicament of Absorption-dominated Reionization: Increased Demands on Ionizing Sources. bibcode 2021ApJ...918L..35D \par [Madau2017] Madau (2017). Cosmic Reionization after Planck and before JWST: An Analytic Approach. bibcode 2017ApJ...851...50M

\end{document}
